% RESEARCH_DRAFT_SHA256: f5c064cf7ff654fc26f4e123a04eda97bfa8517de9be64a3847a5affd8d83799
% Layout and source-reference wrappers expanded for the web renderer.
\documentclass[11pt]{article}
\usepackage[margin=25mm]{geometry}
\usepackage{fontspec,unicode-math}
\setmainfont{Latin Modern Roman}
\setmathfont{Latin Modern Math}
\usepackage{amsmath,amssymb,mathtools,xurl,hyperref}
\hypersetup{colorlinks=true,urlcolor=blue,linkcolor=blue}
\setcounter{secnumdepth}{0}
\setlength{\parindent}{0pt}
\setlength{\parskip}{5pt}
\setlength{\emergencystretch}{3em}
\newcommand{\sref}[2]{\hyperlink{source:#1:#2}{[S#2]}}
\newcommand{\cataloguescope}[1]{\paragraph{Source record.}#1}
\newcommand{\floor}[1]{\left\lfloor#1\right\rfloor}
\begin{document}
\section*{Research-attempt provenance}
This standalone research attempt retains the complete original definition body below,
including its historical source and preparation notes. Those notes describe the
earlier catalogue preparation; the new attempt and its verification are identified
separately after that body. No original ID, catalogue status, or source statement is changed.

% UnsolvedMath ID 20002169; source catalogue code AIM-INFRASTRUCTURE-0069
\section{Certified LaTeX/Markdown import into PreTeXt}\label{20002169}
{\small UnsolvedMath ID 20002169 · Computer Science}
\subsection{Definitions and mathematical statement}

\paragraph{Shared notation.}Standard number systems are $\mathbb N,\mathbb Z,\mathbb Q,\mathbb R,\mathbb C$. Specialized notation and index ranges follow the cited source and the definitions below.


LaTeX, Markdown, and PreTeXt are different document-source representations. Conversion must interpret the supported source constructs and map them into the target structure. Pandoc is a document-conversion tool named in the extract. Ease of conversion can concern supported syntax, user effort, diagnostics, or retained semantics.
\paragraph{Record scope and qualifications.}
The catalogue title is retained as metadata. Named certificates, contracts, bounds or frameworks in that title are not asserted as results of the original source extract.
No comparison metric or input subset is defined; the statement is not a proved performance comparison with Pandoc.

Catalogue screening limitation: PreTeXt software design, documentation, or publication feature request; no mathematical problem is specified by the supplied passage.

The source snapshot labels this record partially solved; this is catalogue metadata rather than independent verification of a mathematical result or the current completion of a project.
The source text below is rendered literally where necessary; damaged formulas and composite extracts are not silently repaired into a new mathematical claim.
The supplied catalogue marks this formulation as an unverified reconstruction; missing assumptions are not certified by the source link.
\paragraph{Statement recorded in the supplied source.}
Converting from latex/markdown easier than pandoc. \href{https://github.com/PreTeXtBook/community/wiki/Moderated-Problem-Session}{[S2]}
\paragraph{Source status.}
The supplied catalogue labels this source record partially solved; that metadata was checked against the supplied snapshot on 2026-10-11. The label is preserved for provenance and is not an independent verification of present open status or mathematical correctness.
\subsection{Short English statement}
Make conversion from LaTeX or Markdown to PreTeXt easier.
\subsection{Sources}
\begin{itemize}
\item[S1] Ulam AI and the UnsolvedMath Contributors, UnsolvedMath: A Curated Collection of Open Mathematics Problems, record 20002169 (AIM-INFRASTRUCTURE-0069). Catalogue curation and metadata: CC BY 4.0. The source extract is retained; definitions, English summaries and qualifications are editorial additions.
\url{https://huggingface.co/datasets/ulamai/UnsolvedMath}.
\item[S2] AIM Workshop Problem Lists, source linked by catalogue record AIM-INFRASTRUCTURE-0069.
\url{https://github.com/PreTeXtBook/community/wiki/Moderated-Problem-Session}.
\end{itemize}
\paragraph{Preparation and review.}
The original source and UnsolvedMath contributors retain attribution for the source extract. Codex (GPT-6; exact serving revision not exposed) assisted with definitions, English summaries, scope notes and formatting. No new human mathematical authorship or independently verified solution is asserted. Permanent source ID 20002169 and code AIM-INFRASTRUCTURE-0069 checked against the complete supplied registry and current repository. Distinct upstream records, including repeated questions and broad or damaged extracts, are retained for complete coverage of the source CS category. A repeated record is not a new independent conjecture. Complete mathematical-equivalence and current-literature review has not been performed. Preparation and final review follow contribution prompts C and D, with source gaps disclosed. Structural and rendering checks do not establish a new theorem or open status.
\label{end:20002169}

\section{Literature and context check}
The primary community-session entry was read. It is an implementation or documentation agenda, and its editorial catalogue title supplies no theorem specification. The following explicitly chosen model tests one route; no actual PreTeXt conversion, documentation release, or user study is claimed.
\section{Research attempt: Balanced structure defeats a finite-state import shortcut}
Consider the restricted syntax of balanced braces, before any full LaTeX or Markdown semantics are added. A proposed importer using a finite-state recognizer with $m$ states cannot recognize all balanced-brace strings exactly. Read the $m+1$ prefixes consisting of $0,1,\ldots,m$ opening braces. Two prefixes of different lengths $i<j$ reach the same state. Append $i$ closing braces to each. The first completed string is balanced, while the second retains $j-i>0$ unmatched openings. Since the recognizer had the same state before the identical suffix, it must make the same decision on both, a contradiction.

A stack parser fixes this particular obstruction. Scan left to right, push an opening brace, and pop for a closing brace; reject a pop from an empty stack and reject a nonempty stack at the end. After every prefix, the stack height equals the number of openings minus closings, unless a negative prefix already caused rejection. This invariant follows by induction on the next character. A string is accepted exactly when no prefix has negative balance and final balance is zero, the usual balanced-brace condition. Each symbol incurs one constant-time stack operation, so the parser takes linear time and at most linear stack space.

This establishes only a structural primitive. Escaping, comments, macro expansion, math syntax, and the differing conventions of source languages remain outside it. A real converter may use richer parsing than a finite-state rule, so the impossibility is not an indictment of all text tools. It pinpoints why an unrestricted nesting promise needs an appropriate memory model, while no amount of merely accepting balanced braces supplies a semantic PreTeXt translation.
\section{Verification and adversarial audit}
The finite examples and implications are checked directly in the argument. No production PreTeXt build or conformance test was run. The proof applies only to the declared model, with its input and identity assumptions.
\paragraph{Audit corrections.} The lower bound is proved for finite-state recognition of the declared brace language. It is not an unsupported impossibility claim about all regular-expression engines or all importers.
\paragraph{External computation.} None; symbolic arguments and boundary checks only.
\section{Scope and remaining gaps}
The result is an auxiliary route analysis, not completion of the original software or documentation agenda. Real versions, supported syntax, transformations, and usability requirements need concrete implementation and testing. No mathematical novelty or current feature-completion claim is made.
\paragraph{Integration review.} Codex (GPT-6; exact serving revision not exposed) reviewed the written arguments and expanded layout and source-reference wrappers. Repository metadata and the production web renderer were checked using Python 3.13 and JavaScript with MathJax 3.2.2. This model-assisted review is not independent human proof verification; no PDF compilation or formal proof checking is asserted.

\section{Completion Estimate}
\noindent\textbf{COMPLETION: 0\%}
The calculations below the preserved source statement investigate a concrete route,
but establish no substantiated new progress toward the entire original target.
Known special cases and failed sufficient conditions are not credited as a solution.
\section{References}
The original source references and catalogue attribution are preserved in the
Sources section above. The following URLs were actually consulted for this attempt;
retention of an additional original reference does not claim fresh consultation of it.
\begin{itemize}
\item \url{https://github.com/PreTeXtBook/community/wiki/Moderated-Problem-Session} (consulted 2026-10-11). 
\end{itemize}
\section{Model and contributor disclosure}
OpenAI gpt-6-astra, requested reasoning configuration ultra, generated this research
attempt on 2026-10-11. The exact serving revision is not exposed. Codex GPT-6
performed orchestration and repository-format review; its exact serving revision
is also not exposed. The mathematical route, boundary checks and final adversarial
self-audit followed contribution prompts A then B in
\texttt{docs/contribution\_prompts.md}, after reading \texttt{CONTRIBUTING.md}.
Python 3.14.2 assembled the standalone documents; any mathematical computation is
specified in Verification. No theorem prover was used. No human mathematical
authorship, substantial human proof correction, or independent human proof checking
is asserted. Model-assisted review, source retrieval, and structural checks do not
establish correctness or novelty. Full private reasoning transcripts are unavailable
as reproducibility artifacts; the complete arguments, inputs and outcomes actually
used are stated here or in the companion batch audit.
\section{Final status}
Unresolved as a new research contribution. This is a concrete route investigation,
not a claim of new complete resolution or a current-status certification. A known
answer reproduced for context does not become a new contribution. The catalogue's
source-status label is preserved; this attempt status does not assert that a
source-labelled solved question is open.
\end{document}
